% Abstract (<= 250 words; count with: detex sections/abstract.tex | wc -w)
\begin{abstract}
Autonomous management of a 5G cell carried by an unmanned aerial vehicle (UAV)
must keep search-and-rescue robots commandable without an operator and without
knowing the terrain. Over an OpenAirInterface and Open5GS testbed serving four
robots, we built and flew a three-loop management stack: a per-robot uplink cap
(every 0.2\,s), a liveness monitor (every 20\,s), and a terrain-blind learning
layer that forecasts link failure and places the UAV in three dimensions (every
30\,s). The decisive loop was the simplest. Long forest outages were not
coverage loss but uplink queuing inside the user equipment (UE): constant-rate
video exceeded what the links drained and held command-and-control (C2) traffic
minutes behind it. The cap, driven only by round-trip time and bytes sent,
raised forest C2 availability from 66.8\,\% to 91.8\,\% in a patrol test and
cut the median C2 outage from 143\,s to 17\,s in a placement run. A second
silent failure, dead UE processes, had contaminated 12.4\,\% of the training
data. With both handled, the learning layer forecasts C2 failure 20\,s ahead
with modest skill that falls as model capacity grows, and neither a linear nor
a gradient-boosted model recovers how terrain changes the effect of altitude.
Over repeated missions with the cap, no learned placement beat fixed hover:
in open field every arm was within 5 points of it, and in forest the arms that
moved the UAV trailed it. For autonomous management
of such cells, sound inner loops, with bounded queues and verified liveness,
must come before learned placement.
\end{abstract}
